
Weight-space encoders map the parameters of a trained neural network to a fixed-size representation from which properties such as test accuracy can be predicted. A feedforward network with hidden layers of widths n₁, …, n_L admits ∏ₗ nₗ! permutation-equivalent weight configurations that are functionally identical but appear as distinct inputs to a flat MLP encoder. At matched capacity, this creates a representational burden: encoder parameters are consumed by orbit disambiguation rather than accuracy-predictive feature learning.

Permutation-equivariant encoders such as Neural Functional Networks (NFNs) address this by construction: all permutation-equivalent configurations are mapped to the same embedding. The empirical magnitude of this advantage, however, has not been measured under capacity-controlled conditions. Prior comparisons between NFN and flat MLP encoders do not match encoder parameter counts, so the observed performance differences could reflect both inductive bias advantages and capacity differences.

\textbf{Research gap.} No published study reports a Δρ benchmark comparing matched-capacity (±5\%) NFN and flat MLP encoders on the Schurholt model zoo benchmarks with bootstrap 95\% confidence intervals, an intermediate invariant baseline (Deep Sets), and direct mechanism measurement via permutation sensitivity scores.

\textbf{This work.} We train flat MLP, Deep Sets, and NFN encoders at matched approximately 500K parameters on the Schurholt MNIST-CNN model zoo (hyp_rand variant) and report:

\begin{enumerate}
\item Permutation orbit existence (H-E1): whether the zoo contains non-trivial permutation orbits.
\item Flat MLP sensitivity (H-M1): whether flat MLPs exhibit high permutation sensitivity at this capacity.
\item NFN sensitivity (H-M2): whether NFN encoders achieve near-zero permutation sensitivity.
\item Primary performance comparison (H-M3): Δρ with paired bootstrap 95\% CI.
\end{enumerate}

The paper is organized as follows. Section 2 reviews related work. Section 3 describes the methodology. Section 4 presents the experimental setup. Section 5 reports results. Section 6 discusses findings and limitations. Section 7 concludes.

